\clearpage
\fancyhead[L]{\small GAUSS AN BODE.}
\hypertarget{native-438-part-lower}{}
\begin{center}
{\large GAUSS AN BODE.}

Astronomisches Jahrbuch für 1823, S. 228..230. Berlin 1820.
\end{center}

\hfill Göttingen, den 3. September 1820.

\par\indent Wie gewünscht übersende ich Ihnen hier die Ephemeride für den Lauf der \emph{Pallas} während ihrer Sichtbarkeit im Jahr 1821. Sie wurde diesmal von Herrn \textsc{von Staudt} berechnet, einem jungen Mann von ausgezeichnetem Talent, der sich hier dem Studium der Mathematik und Astronomie widmet. Die Berechnung ist erst heute fertig geworden, und Sie wollten sie so bald wie möglich erhalten. Deshalb kann ich die noch nicht abgeschlossene genaue Vorausberechnung der Opposition vom 19. Mai 1821 nicht beifügen. Sonst könnte ich die Sendung erst einen Posttag später abschicken.

Den Meridiankreis von \textsc{Reichenbach} habe ich in unseren gelehrten Anzeigen ausführlich beschrieben. Es wird daher nicht nötig sein, hier etwas davon zu wiederholen. Ich habe seit dem 21. Februar wohl etwa 1200 Beobachtungen damit gemacht, von denen bisher allerdings erst ein Teil vollständig reduziert ist. Diesmal teile ich Ihnen die bis zum 15. Juli beobachteten Sonnendeklinationen mit (von den gleichzeitig am Mittagsfernrohr beobachteten Rektaszensionen ist erst ein Teil reduziert). Der Ort des Pols auf dem Instrument ist dabei bereits anhand sämtlicher dafür geeigneter Beobachtungen so genau wie möglich bestimmt. Die Beobachtungen sind um die Refraktion nach \textsc{Bessels} Tafel und auch um die Parallaxe bereinigt. Für die mittlere Horizontalparallaxe wurde $=8''60$ angesetzt. Die Flexion blieb jedoch unberücksichtigt. Nach meinen bisherigen Versuchen scheint sie fast ganz unmerklich zu sein. Sollte sie aber doch vielleicht \emph{einige Zehntel einer Sekunde} betragen, so lässt sich diese Größe erst durch eine lange Reihe von Beobachtungen ermitteln.


\Needspace{17\baselineskip}
\SourceLeaf{439}
\begin{center}
\emph{Beobachtete Deklinationen der Sonne} 1820.

O oberer Rand \quad U unterer Rand
\end{center}

\begin{center}
\begin{tabular}{ll|r r r l}
\multicolumn{6}{c}{Südlich} \\
Febr. & 22 & $10^\circ$ & $12'$ & $36''5$ & O \\
 & 27 & $8$ & $54$ & $23.0$ & U \\
 & 28 & $7$ & $59$ & $29.9$ & O \\
 & 29 & $8$ & $9$ & $11.3$ & U \\
März & 5 & $6$ & $14$ & $26.5$ & U \\
 & 11 & $3$ & $53$ & $5.3$ & U \\
 & 14 & $2$ & $11$ & $0.0$ & O \\
 & 19 & $0$ & $44$ & $37.2$ & U \\
\multicolumn{6}{c}{Nördlich} \\
April & 3 & $5$ & $38$ & $52.1$ & O \\
 & 4 & $6$ & $1$ & $48.2$ & O \\
\end{tabular}
\end{center}

\Needspace{15\baselineskip}
\begin{center}
\begin{tabular}{ll|r r r l}
\multicolumn{6}{c}{Nördlich} \\
April & 12 & $9^\circ$ & $0'$ & $52''2$ & O \\
 & 13 & $8$ & $50$ & $43.1$ & U \\
 & 16 & $9$ & $55$ & $14.7$ & U \\
 & 18 & $11$ & $9$ & $15.3$ & O \\
 & 21 & $12$ & $11$ & $13.2$ & O \\
 & 22 & $11$ & $59$ & $26.2$ & U \\
 & 23 & $12$ & $19$ & $29.7$ & O \\
 & 24 & $13$ & $11$ & $9.9$ & O \\
 & 26 & $13$ & $50$ & $10.3$ & O \\
 & 27 & $13$ & $37$ & $30.7$ & U \\
Mai & 6 & $16$ & $19$ & $13.5$ & U \\
\end{tabular}
\end{center}

\Needspace{15\baselineskip}
\begin{center}
\begin{tabular}{ll|r r r l}
\multicolumn{6}{c}{Nördlich} \\
Mai & 11 & $18^\circ$ & $11'$ & $25''9$ & O \\
 & 19 & $20$ & $4$ & $23.5$ & O \\
 & 27 & $21$ & $4$ & $11.3$ & U \\
 & 31 & $21$ & $41$ & $11.8$ & U \\
Juni & 6 & $22$ & $25$ & $19.1$ & U \\
 & 20 & $23$ & $43$ & $27.5$ & O \\
 & 25 & $23$ & $50$ & $30.6$ & O \\
 & 27 & $23$ & $4$ & $58.1$ & U \\
Juli & 11 & $22$ & $23$ & $0.0$ & O \\
 & 12 & $22$ & $14$ & $51.1$ & O \\
 & 15 & $21$ & $16$ & $29.3$ & U \\
\end{tabular}
\end{center}

Ich teile Ihnen noch die Ergebnisse meiner Beobachtungen von 20 Zenitsternen mit, das heißt die auf den Anfang von 1820 reduzierten mittleren Zenitdistanzen. Für die Reduktion wurden dieselben Elemente verwendet, die \textsc{Bessels} Tafeln für den Polarstern zugrunde liegen. Es sind dieselben Sterne, die Professor \textsc{Schumacher} mit dem Zenitsektor von \textsc{Ramsden} 1819 in Lauenburg und in diesem Jahr in Skagen beobachtet hat. Aus diesen Beobachtungen wird die Krümmung der verschiedenen Abschnitte des Meridianbogens von Skagen bis Göttingen abgeleitet werden. Ich kann diese Krümmungen noch nicht angeben, da mir die \emph{Endergebnisse} der Beobachtungen von \textsc{Schumacher} noch nicht bekannt sind. Ich bemerke nur noch, dass ich jeden Stern \emph{mindestens} dreimal in der einen und dreimal in der anderen Lage des Kreises beobachtet habe. Bei der Reduktion auf den Anfang von 1820 wurde die Eigenbewegung der Sterne nicht berücksichtigt. Sie lässt sich aber mit größter Genauigkeit aus dem Vergleich meiner Bestimmungen mit denen von \textsc{Brandley} ableiten. Dabei ist die Polhöhe der hiesigen Sternwarte mit $=51^\circ\,31'\,48''7$ anzusetzen. Das Mittel meiner Beobachtungszeiten fällt in die ersten Augusttage.

\Needspace{27\baselineskip}
\begin{center}
\begin{tabular}{l|l|l}
Sterne. & M. Z. D. 1820. & Kollimationsfehler \\
\hline
$\eta$ Cephei & $9^\circ\,36'\,44''28$ N. & $2'\,15''50$ \\
\emph{Piazzi XX.} 222. & $8\;\;17\;\;7.34$ & $16.65$ \\
$\iota$ Draconis & $8\;\;4\;\;11.09$ & $16.78$ \\
47 Draconis & $7\;\;38\;\;26.35$ & $17.19$ \\
Cephei 2 Hev. & $6\;\;50\;\;16.06$ & $16.96$ \\
48 Draconis & $6\;\;2\;\;55.95$ & $17.36$ \\
53 Draconis & $5\;\;1\;\;32.34$ & $17.06$ \\
33 Cygni & $4\;\;29\;\;25.71$ & $18.35$ \\
49 Draconis & $3\;\;52\;\;25.79$ & $17.18$ \\
46 Draconis & $3\;\;49\;\;49.51$ & $16.62$ \\
\emph{Piazzi XX.} 391. & $2\;\;17\;\;59.29$ & $17.64$ \\
51 Draconis & $1\;\;35\;\;38.56$ & $16.96$ \\
$\kappa$ Cygni & $1\;\;30\;\;37.18$ & $17.20$ \\
\emph{Piazzi XXI.} 32. & $1\;\;18\;\;6.93$ & $16.27$ \\
20 Cygni & $1\;\;0\;\;18.21$ & $17.54$ \\
$\iota$ Cygni & $0\;\;10\;\;46.28$ S. & $17.06$ \\
$c'$ Cygni Praec. & $1\;\;25\;\;6.27$ & $18.03$ \\
$c'$ Cygni Sequ. & $1\;\;25\;\;33.33$ & $16.37$ \\
$\theta$ Cygni & $1\;\;43\;\;15.97$ & $17.79$ \\
$\omega'$ Cygni & $2\;\;44\;\;22.59$ & $17.71$ \\
\multicolumn{2}{r|}{Mittlerer Kollimationsfehler . .} & $2'\,17''22$ \\
\end{tabular}
\end{center}

Sie sehen, dass die Bestimmung des Kollimationsfehlers nur bei zwei Sternen um mehr als $1''$ vom Mittelwert abweicht.
